\section{Musical Concepts Used by the Language}\label{sec:music-theory-primer}

To understand the abstractions Motivo provides, we need a shared vocabulary for western tonal music.
The language organizes material both vertically (pitch and harmony) and horizontally (time and rhythm).
This section is not a complete theory course; it covers only the concepts referenced explicitly in later
chapters.

\subsection{Pitch, Notes, and Octaves}\label{detail-subsec:pitch-notes-and-octaves}

The fundamental unit of melody and harmony is the \textit{note}, which represents a sound of a specific frequency.
In western music, the octave is divided into twelve semitones.
The seven natural notes are designated by the letters A through G\@.
\textit{Accidentals} modify these natural notes: a sharp (\#) raises the pitch by one semitone, while a flat (b) lowers
it by one semitone.

An \textit{octave} is the interval between one musical pitch and another with double its frequency.
To uniquely identify a note across the entire audible spectrum, it is denoted by its pitch class and its octave number
(e.g., \texttt{C4}, which is often referred to as ``Middle C'').
When two or more distinct notes are sounded simultaneously, they form a \textit{chord}, which provides the harmonic
foundation of a composition.

\subsection{Time, Rhythm, and Tempo}\label{detail-subsec:time-rhythm-and-tempo}

While pitch defines the frequency of a sound, rhythm defines its placement and duration in time.
The speed at which a piece of music is played is governed by its \textit{tempo}, typically measured in Beats Per Minute
(BPM).
A tempo of 120 BPM signifies that there are 120 quarter-note beats in one minute, or one beat every 0.5 seconds.

Music is structurally divided into \textit{measures} (or bars) by a \textit{time signature}.
A time signature, such as 4/4, consists of two numbers: the numerator indicates the number of beats per measure (four),
and the denominator indicates the note value that constitutes one beat (the quarter note).
A \textit{rest} is a period of silence that has a specific duration, functioning as a ``silent note'' that preserves the
rhythmic alignment of the timeline.

\subsection{Textual Notation and ABC}\label{subsec:abc-notation}

Long before DAWs, composers captured music in written notation.
Modern interchange formats inherit that declarative tradition in lighter-weight textual forms.

\topic{ABC notation.} Is a plain-text format for notating folk and traditional music~\cite{Walshaw_2011}.
Header fields encode metadata (title, meter, key); the body encodes notes, durations, and repeats using an ASCII
alphabet.
ABC is human-readable and widely supported by conversion tools, but it is a \textit{notation} format rather than a
programming language: it lacks parameterized patterns, typed variables, and imperative control flow.
MusicXML addresses score interchange between notation programs through a structured XML representation~\cite{Good_2001}.

We mention ABC because it illustrates a useful contrast for Motivo.
Notation formats excel at fixed scores intended for human performance; Motivo instead targets procedural descriptions
that compile into event streams.
Neither ABC nor MusicXML provides the executable abstractions---typed patterns, loops, and conditional
generation---that Motivo offers at the source level.
